\documentclass[book.tex]{subfiles}

\begin{document}
\chapter{Geometry Processing}
\label{chap:geometry_processing}
\noindent\rule{11cm}{0.4pt} \\
\textbf{Prerequisites:}  \autoref{chap:ml_basics}, \autoref{chap:probability}  \\
\textbf{Difficulty Level:} ** \\
\noindent\rule{11cm}{0.4pt}
\vspace{5mm}
\textcolor{red}{TODO: This is  sciml talk notes from talk by Yu Wang}
The shape analysis community in graps has a long history of building group transformation invariant algorithms.

\section{Shape Analysis}

\textcolor{red}{TODO: Look at shape analysis review article with Yu Wang and Justin Solomon}

Two tasks
\begin{itemize}
    \item Distance
    \item Segmentation
    \item Shape Description
    \item Correspondance: Given two shapes, find 1-1 mapping.
\end{itemize}

Shape analysis is the generalization of image analysis to 3D. It's something of a 2.5D problem since the boundary of the image is what ends up really matters.


\section{Operator Approach to Shape Analysis}

For example, consider a mesh representation of a shape. There are many possible meshes you can use. 

We define a shape as the discretization of a continuous operator on a surface (\textcolor{red}{TODO: I'm not entirely sure how this works.}) You use the Laplacian operator

\begin{align}
    \mathbb{R}^2: \Delta^2: \frac{\partial^2}{\partial x^2} + \frac{\partial^2}{\partial y^2}
\end{align}

The Laplace-Beltrami operator is a linear operator defined on a manifold. You can use the finite element method to discretize the operator. This will result in the derivation of a "cotangent" Laplacian.

\begin{align}
    \Delta_M: \frac{1}{\sqrt{|det g|} \partial_i(\sqrt{|det g|})g^{ij}...
\end{align}

Some you have some discretized geometry operator $S: \mathbb{R}^{n \times n}$. This type of analysis is called spectral analysis.

\textcolor{red}{TODO: How does this relate to the rest of the book?}

The goal here is to create an innvariance to the representation of the mesh.

ShapeDNA advocates using Laplacian eigenvalues as vector representation for shape retrieval.

Operator spectrum provides a shape2vector approach. Shape -> eigenvalues -> semantics. Map from eigenvalues to semantics can be learned.

Why does a operator representation have robustness to noise? The lower eigenvalues are more stable to perturbation

Heat kernel signatures. Here $k_t$ is also known as a Green's funciton.

\begin{align}
    h_t(x) &= k_t(x, x) \\
    k_t(x, y) &= \sum_{i=0}^\infty e^{-\lambda_i t} \phi_i(x) \phi_j(y) 
\end{align}

$h_t(x)$ is a function for each point on the surface. 

Wave kernel signature contains a filter signature other than $f(\lambda) = e^{-\lambda t}$. Optimal spectral decompositions consider a differetn filter.
\section{Intrinsic and Extrinsic Operators}

Extrinsic geometry care about the embedding of a shape in $\mathbb{R}^3$. Intrinsic shapes are invariant to isometry. An isometry is length preserving. The Laplacian operator is invariant to isometries.

An intrinsic geometry is any origami equivalent to a flat piece of paper.

People usually work with intrinsic operators but extrinsic operators are starting to become popular.


\section{Operators for Geometric Analysis}
\subsection{Dirichlet to Neumann Operators $\mathcal{S}$}

Consider a volume $\Omega$ bounded $\Gamma = \partial \Omega$. The neumann data is 
\begin{align}
    g_n &= \frac{\partial}{\partial n} u(\Gamma)
\end{align}
DtN operator is
\begin{align}
    g \mapsto g_n
\end{align}
Also called the Stekhov-Poincare operator. $\mathcal{S}$ can be written as a composition of 4 operators, boundary operators. This can be generalized to open surface.

The operator $\mathcal{V}$ the single layer potential is defined as
\begin{align}
    [V\phi](x) &= \int_\Gamma G(x, y) \phi(y) dy
\end{align}
\begin{align}
    v(x, y) &= \frac{1}{|x-y|}
\end{align}
is an inverse distance operator. The operator $\mathcal{S}$ encodes the extrinsic geometry of the system.

Level sets of Steklov eigenfunctions conform to mean curvatures.

\subsection{Dirac Operators}

\begin{align}
    D = \sqrt{\Delta}
\end{align}
\textcolor{red}{TODO: Look at "Surface Networks" paper. Also look up extrinsic dirac operators paper}

\section{Learnable Operators}

\textcolor{red}{TODO: Look up HodgeNet, DiffusionNet, PD-MeshNet, MeshCNN and others}

\end{document}